\section{Introduction}

Believe it or not I am culminating four years of mathematics with a project about soap films - those elegant shapes that captivate the attention and imagination of every 2 year old (and for that matter a great number of adults)! Behind the beauty lie complicated mathematical problems that have been propelling forward numerous fields of mathematics since the 18th century when Jean Baptiste Meusnier (1754-1793) discovered the first examples other than the plane, in the Catenoid and Helicoid which we'll meet again later.

What makes these soap film surfaces special is that, given an arbitrarily complicated boundary, soap films minimize surface area (or in fact more generally they form at turning points of surface area). As will become apparent later all minimal surfaces satisfy the condition that they have zero mean curvature.

In the mid-nineteenth century Joseph Plateau carried out experiments accurately measuring the properties of surfaces formed on wire frames. His results led him to conjecture that every closed curve which doesn't intersect itself can be spanned by a minimal surface. The problem of finding a minimal surface for a given boundary curve became known as \emph{Plateau's Problem}.

It would appear that soap films themselves answer this question but to formulate a mathematical treatment has proved more difficult.

Lagrange was able to formulate the conditions for a minimal surface into a partial differential equation but lacked the tools to solve it in order to prove the conjecture. This has spurred on developments in partial differential equations, calculus of variation, complex analysis, and differential geometry.

A proof didn't appear for \emph{Plateau's Problem} until 1931 when Jesse \nohyphens{Douglas} and independently Tibor Rado, proved existence for a given boundary curve. This was a high point for the calculus of variations but said very little about the geometric properties of the surfaces or in fact how to find them! Therefore, as with so much in mathematics, the answers provided a wealth of new problems.

Research into minimal surfaces over the past 30 years has concentrated on the more open question of minimal surfaces which lack boundary curves. Further more armed with the power of Riemannian Geometry, has generalized them into higher dimensional minimal manifolds.

So is this entire field \emph {mathematics for the sake of mathematics} with the only connection to the real world being soap bubbles? No! Minimal surfaces have numerous real world uses. To understand why, we need to note some of the properties of these surfaces. Minimal surfaces: 
\begin{itemize}
	\item Minimize surface area.
	\item Form in soap films because their surface tension is in equilibrium at each point in the surface.
\end{itemize}

The most obvious real world example of minimal surfaces is in architecture - notably in the construction of roofs, having a minimal surface area, minimizes weight and cost of materials. Further by having surface tension in equilibrium at each point in a roof stabilizes the whole structure making it stronger. Figure \ref{fig:munich} shows a section of the roof of the Olympic Stadium in Munich the entirety of which is a minimal surface. As you can see while the properties of minimal surfaces make them ideal for such purposes they are also very aesthetically pleasing!

\begin{figure}[htbp]
	\centering
       \includegraphics[width=8cm]{Images/Munich.eps}
   \caption{Munich Olympic Stadium \cite{MUN}}
   \label{fig:munich}
\end{figure} 


They also crop up widely in molecular engineering, crystallography and in the future, nanotechnology.